% -*- coding: utf-8 -*-
% !TEX root = ../main.tex

La figure~\ref{fig:synthese_pipeline} résume la chaîne d'inférence principale.
Le modèle Word2Vec est préentraîné avec l'objectif CBOW, puis figé. Lors de
la classification, seule sa table d'embeddings d'entrée est utilisée pour
associer un vecteur à chaque unité sous-lexicale. L'agrégateur transforme
ensuite la séquence de longueur variable en un vecteur de taille fixe, qui est
fourni à la tête de classification.

\begin{figure}[!ht]
    \centering
    \resizebox{\textwidth}{!}{%
        \begin{tikzpicture}[
            node distance=7mm,
            flow/.style={-{Latex[length=2.8mm]}, very thick, draw=black!70},
            stage/.style={
                draw=black!65,
                rounded corners=2mm,
                minimum height=24mm,
                align=center,
                font=\small,
                inner xsep=4mm,
                inner ysep=3mm
            },
            input/.style={stage, fill=gray!12, minimum width=22mm},
            representation/.style={stage, fill=blue!9, minimum width=36mm},
            aggregation/.style={stage, fill=orange!13, minimum width=38mm},
            classifier/.style={stage, fill=green!11, minimum width=43mm},
            output/.style={stage, fill=red!9, minimum width=35mm},
            annotation/.style={
                align=center,
                font=\scriptsize,
                text=black!75,
                inner sep=1mm
            }
        ]

        \node[input] (text) {
            \textbf{Texte}\\[1mm]
            $t$
        };

        \node[representation, right=of text] (bpe) {
            \textbf{Tokenisation BPE}\\[1mm]
            $(u_1,\ldots,u_n)$\\
            $u_j\in\mathcal{V}$
        };

        \node[representation, right=of bpe] (cbow) {
            \textbf{Embeddings CBOW}\\[1mm]
            $E\in\mathbb{R}^{|\mathcal{V}|\times d}$\\
            $\boldsymbol e_j=E[u_j]\in\mathbb{R}^{d}$
        };

        \node[aggregation, right=of cbow] (mean) {
            \textbf{Agrégateur moyen}\\[1mm]
            $\displaystyle
             \boldsymbol x=\frac{1}{n}\sum_{j=1}^{n}\boldsymbol e_j$\\
            $\boldsymbol x\in\mathbb{R}^{d}$
        };

        \node[classifier, right=of mean] (head) {
            \textbf{Tête de classification}\\[1mm]
            Linéaire $d\rightarrow h_1$ + ReLU\\
            Linéaire $h_1\rightarrow h_2$ + ReLU\\
            Linéaire $h_2\rightarrow1$
        };

        \node[output, right=of head] (decision) {
            \textbf{Sortie binaire}\\[1mm]
            $z\in\mathbb{R}$ : un logit\\
            $\widehat p=\sigma(z)\in[0,1]$\\
            $\widehat y=\mathbf{1}[\widehat p\geq0{,}5]\in\{0,1\}$
        };

        \draw[flow] (text) -- (bpe);
        \draw[flow] (bpe) -- node[annotation, above] {$n$ unités} (cbow);
        \draw[flow] (cbow) -- node[annotation, above]
            {$n\times d$ valeurs} (mean);
        \draw[flow] (mean) -- node[annotation, above]
            {$\boldsymbol x\in\mathbb{R}^{d}$} (head);
        \draw[flow] (head) -- node[annotation, above]
            {un neurone de sortie} (decision);

        \node[annotation, above=5mm of cbow] (offline) {
            préentraînement CBOW hors ligne, puis paramètres figés
        };
        \draw[-{Latex[length=2mm]}, dashed, draw=black!55]
            (offline) -- (cbow);

        \end{tikzpicture}%
    }
    \caption{Synthèse du pipeline d'inférence pour la classification binaire.}
    \label{fig:synthese_pipeline}
\end{figure}

Le vocabulaire BPE est noté $\mathcal{V}$, $n$ désigne le nombre d'unités du
texte et $d$ la dimension d'un embedding. L'agrégation moyenne conserve cette
dimension indépendamment de $n$ : elle transforme donc une matrice de
$n\times d$ valeurs en un unique vecteur $\boldsymbol{x}\in\mathbb{R}^{d}$.
Les symboles $h_1$ et $h_2$ désignent les largeurs des deux couches cachées de
la tête. Leurs valeurs concrètes seront introduites dans la formulation du
problème.

La dernière couche ne produit pas deux scores, mais un unique logit
$z\in\mathbb{R}$. Après application de la sigmoïde, ce logit devient la
probabilité estimée de la classe positive. L'étiquette réelle est
$y\in\{0,1\}$ et la prédiction associée est
$\widehat y\in\{0,1\}$.
